\documentclass[11pt,a4paper]{article}
\usepackage[margin=1in]{geometry}
\usepackage{amsmath}
\usepackage{booktabs}
\usepackage{graphicx}
\usepackage{caption}
\usepackage[hidelinks]{hyperref}
\usepackage{parskip}
\usepackage{times}

\captionsetup[table]{font=small,labelfont=bf,skip=4pt}

\title{\textbf{The Variance Risk Premium and a Defined-Risk Option Strategy}\\[4pt]
\large Evidence from S\&P 500 options, 1990--2026}

\author{Alex Yuan\\[2pt]\small Independent research project}
\date{August 2026}

\begin{document}
\maketitle

\begin{abstract}
\noindent
This study tests whether the variance risk premium (VRP) --- the tendency for option-implied
volatility to exceed subsequently realised volatility --- exists, persists, and is captured by
a defined-risk option strategy. Using 9,178 daily observations of the CBOE Volatility Index (VIX) against
subsequent 21-day realised volatility of the S\&P 500 from 1990 to 2026, the mean premium is
\textbf{+4.10 volatility points} (HAC 95\% CI [3.57, 4.62]). The result is robust to autocorrelation
correction, distribution-free bootstrap, and decade-by-decade subsampling. A na\"ive t-test overstates
the evidence by a factor of four, illustrating the consequences of ignoring overlapping windows.
A simulated bull put spread capturing this premium returns \textbf{\$8.35 per trade} net of \$3.50
costs across 9,201 trades, with no decay out-of-sample. Nine volatility-based timing filters were
tested; \textbf{all nine underperformed unconditional entry}. The premium is shown to be compensation
for crash risk rather than a market inefficiency: the distribution has skewness of $-3.65$ and every
losing year coincides with an equity bear market. Principal limitations are the omission of volatility
skew and the assumption of theoretical fills.
\end{abstract}

\vspace{4pt}
\noindent\fbox{\begin{minipage}{0.97\textwidth}
\textbf{\large Summary of findings}
\smallskip

\begin{tabular}{@{}p{0.52\textwidth}p{0.40\textwidth}@{}}
\textbf{Does the premium exist?} & Yes. +4.10 vol points, $z=15.3$. \\[3pt]
\textbf{Has it decayed since 1990?} & No. Every decade significant; 2020s mean 3.79. \\[3pt]
\textbf{Is it an inefficiency?} & No. Skew $-3.65$; it pays for crash risk. \\[3pt]
\textbf{Does a spread capture it?} & Yes in simulation. \$8.35/trade net. \\[3pt]
\textbf{Does timing by VIX help?} & No. All nine filters were worse. \\[3pt]
\textbf{Is it tradeable after costs?} & Not established. Skew and fills untested. \\
\end{tabular}
\end{minipage}}

\section{Introduction}

Selling options is widely claimed to be profitable on the grounds that ``implied volatility exceeds
realised volatility.'' The claim is rarely accompanied by evidence, and where evidence is offered it
is frequently invalidated by a statistical error: daily observations of a forward-looking volatility
window overlap heavily, violating the independence assumption of standard inference.

This study asks three questions in sequence. First, does the variance risk premium exist in S\&P 500
index options, measured correctly? Second, has it persisted? Third, does a defined-risk option
structure --- the bull put spread --- capture enough of it to matter after costs?

The motivating context is a retail account of approximately £100 --- small enough that per-trade risk
must be tightly bounded.
This constrains the strategy space: cash-secured puts and covered calls require substantial collateral,
whereas narrow vertical spreads bound maximum loss to tens of dollars per position. Whether such a
structure retains any edge after commissions and bid--ask spread is the practical question behind the
analysis.

The paper finds that the premium is large and statistically robust once overlapping windows are handled
correctly; that it has not decayed across three and a half decades, though its risk-adjusted
attractiveness has; that a simulated defined-risk spread captures it profitably under every cost
assumption tested; and that every attempt to time entry by volatility regime made results worse.
Falsification criteria were pre-registered before any data was examined, and are reported against
explicitly in Sections 3 and 4.

\section{Data and Method}

Daily closing values of \texttt{\^{}VIX}, \texttt{\^{}GSPC} and \texttt{\^{}IRX} were obtained from
Yahoo Finance for 2 January 1990 to 14 August 2026, giving 9,178 complete observations. The sample
begins in 1990 because CBOE back-computed the current VIX methodology only to that date.

For each date $t$, realised volatility over the \emph{following} 21 trading days (approximately 30
calendar days, matching the VIX horizon) was computed from log returns as
\begin{equation}
RV_{t \rightarrow t+21} = 100 \times \sqrt{\frac{252}{21}\sum_{i=1}^{21} r_{t+i}^{2}},
\qquad r_t = \ln\!\left(\frac{P_t}{P_{t-1}}\right),
\end{equation}
using the zero-mean convention standard in the realised volatility literature. The variance risk
premium is
\begin{equation}
VRP_t = VIX_t - RV_{t \rightarrow t+21}.
\end{equation}

Window alignment was verified by reconstructing individual observations manually and comparing them
against the vectorised computation. This check detected one substantive error during development ---
a backward-looking window --- which would have produced plausible but meaningless results.

\subsection{Inference under overlapping windows}

Consecutive observations share 20 of their 21 forward returns, inducing serial correlation to lag 21.
Ordinary standard errors are therefore invalid. Three corrections were applied: (i) non-overlapping
subsampling at every 21st observation, repeated across all 21 possible starting offsets; (ii)
Newey--West heteroskedasticity- and autocorrelation-consistent standard errors with \texttt{maxlags}
$=21$; and (iii) a bootstrap with 10,000 resamples drawn from the non-overlapping subsample, which
assumes nothing about distributional shape.

\section{Results I: The Variance Risk Premium}

Mean VIX over the sample was 19.44 against mean subsequent realised volatility of 15.36. The premium
averaged $+4.10$ volatility points, with a median of $+4.71$ and standard deviation of 6.63. It was
positive in 85.8\% of observations.

\subsection{Inference}

\begin{table}[h]
\centering
\caption{Inference on the mean variance risk premium ($n=9{,}178$)}
\begin{tabular}{@{}lccl@{}}
\toprule
Method & $t$ / $z$ & 95\% CI & Assumption \\
\midrule
Na\"ive one-sample $t$-test & 59.2 & --- & Independence (violated) \\
Non-overlapping subsamples & 12.0 -- 13.8 & --- & None; discards 95\% of data \\
Newey--West HAC & 15.3 & [3.57, 4.62] & Correct lag specification \\
Bootstrap (10,000 draws) & --- & [3.31, 4.55] & None \\
\bottomrule
\end{tabular}
\end{table}

The na\"ive test overstates the evidence by approximately a factor of four, closely matching the
$\sqrt{21}\approx 4.6$ predicted from the overlap structure. The regression's Durbin--Watson statistic
of 0.094 (2.0 indicating no autocorrelation) independently confirms severe serial dependence. All three
valid methods agree, and the distribution-free bootstrap brackets the HAC interval, so the conclusion
is not an artefact of any single method's assumptions.

\subsection{Distribution}

\begin{table}[h]
\centering
\caption{Distribution of the variance risk premium (volatility points)}
\begin{tabular}{@{}lrlr@{}}
\toprule
Statistic & Value & Statistic & Value \\
\midrule
Minimum & $-71.96$ & 75th percentile & $+7.23$ \\
1st percentile & $-19.95$ & 95th percentile & $+11.71$ \\
5th percentile & $-4.87$ & 99th percentile & $+16.65$ \\
25th percentile & $+2.21$ & Maximum & $+30.77$ \\
Median & $+4.71$ & Skewness & $-3.65$ \\
Mean & $+4.10$ & Excess kurtosis & $29.27$ \\
\bottomrule
\end{tabular}
\end{table}

The interquartile range, $[+2.21, +7.23]$, is entirely positive: the typical period pays a modest,
unremarkable premium. The asymmetry lies wholly in the tails. The worst observation ($-71.96$) is more
than twice the magnitude of the best ($+30.77$), and the 1st percentile is worse than the 99th
percentile is good.

This is the payoff profile of an insurance seller. The economic interpretation follows directly: the
premium is \emph{compensation for bearing crash risk}, not evidence of mispricing. Table 3 makes the
point concrete.

\begin{table}[h]
\centering
\caption{Six worst episodes (monthly minima)}
\begin{tabular}{@{}lrrrl@{}}
\toprule
Date & VIX & Realised vol & VRP & Event \\
\midrule
2020-02-20 & 15.56 & 87.52 & $-71.96$ & COVID-19 crash begins \\
2020-03-04 & 31.99 & 93.43 & $-61.44$ & COVID-19 crash \\
2008-09-16 & 30.30 & 80.21 & $-49.91$ & Lehman Brothers collapse \\
2008-11-04 & 47.73 & 74.80 & $-27.07$ & Post-Lehman crisis \\
2025-04-02 & 21.51 & 47.90 & $-26.39$ & Tariff shock \\
2011-08-01 & 23.66 & 48.43 & $-24.77$ & US sovereign downgrade \\
\bottomrule
\end{tabular}
\end{table}

The observation of 20 February 2020 is the most instructive in the dataset. VIX stood at 15.56 --- an
unremarkable level implying a calm month ahead --- while realised volatility over the subsequent 21
days was 87.52. The market's own forecast was wrong by a factor of five with no warning priced in.
By contrast, the entry of 4 November 2008 at VIX 47.73 produced a smaller loss: by then the crisis was
visible and priced. \emph{The danger is not high measured volatility; it is low volatility preceding
an unanticipated shock.}

\subsection{Sub-period stability}

\begin{table}[h]
\centering
\caption{Variance risk premium by decade}
\begin{tabular}{@{}lrrrrrc@{}}
\toprule
Period & $n$ & Mean & SD & \% positive & Mean/SD & HAC 95\% CI \\
\midrule
1990s & 2,527 & 5.45 & 4.31 & 92.5 & 1.26 & [4.84, 6.06] \\
2000s & 2,515 & 3.35 & 7.08 & 81.5 & 0.47 & [2.23, 4.48] \\
2010s & 2,516 & 3.68 & 5.44 & 84.3 & 0.68 & [2.92, 4.44] \\
2020s & 1,620 & 3.79 & 9.61 & 84.4 & 0.39 & [1.97, 5.60] \\
\bottomrule
\end{tabular}
\end{table}

Every decade's interval excludes zero, and the 2020s mean exceeds both the 2000s and 2010s: the premium
has not decayed in magnitude. Two qualifications apply. The 1990s interval only just fails to overlap
the 2000s, suggesting a single downward level shift after that decade rather than continuous decay,
with the three subsequent decades statistically indistinguishable. More importantly, dispersion has
more than doubled, so the ratio of premium to variability has fallen by roughly two-thirds (1.26 to
0.39). \emph{The premium is the same size; the risk borne to collect it is substantially greater.}

\subsection{Conditioning on the volatility regime}

\begin{table}[h]
\centering
\caption{Variance risk premium by VIX level}
\begin{tabular}{@{}lrrrrrr@{}}
\toprule
VIX bucket & $n$ & Mean & Median & SD & \% positive & VRP/VIX \\
\midrule
$(0, 15]$ & 2,951 & 2.82 & 3.61 & 4.59 & 86.3 & 0.220 \\
$(15, 20]$ & 2,804 & 4.11 & 5.02 & 5.37 & 87.1 & 0.237 \\
$(20, 25]$ & 1,825 & 4.55 & 5.70 & 6.58 & 83.6 & 0.204 \\
$(25, 30]$ & 864 & 5.56 & 6.83 & 7.56 & 84.4 & 0.204 \\
$(30, 100]$ & 734 & 6.34 & 8.67 & 12.87 & 86.1 & 0.177 \\
\bottomrule
\end{tabular}
\end{table}

Mean premium rises monotonically with VIX, contradicting the intuition that elevated volatility is
dangerous to sell into. However, expressed as a proportion of implied volatility the premium is flat
to slightly declining --- the larger absolute figures at high VIX reflect scaling rather than richer
compensation. On a risk-adjusted basis the highest bucket is the \emph{worst}: mean/SD of 0.49 against
0.77 for the $(15,20]$ bucket. Every bucket contains a catastrophic outcome, so no volatility regime
offers safety.

\subsection{Verdicts against pre-registered criteria}

\textbf{F1} (``mean not significantly positive after correction'') --- \textbf{not triggered}.
\textbf{F2} (``recent sub-period insignificant'') --- \textbf{not triggered}.
\textbf{F3} (``implied per-trade edge below transaction costs'') --- \textbf{unresolved} with index
data; addressed in Section 4.

\section{Results II: A Simulated Defined-Risk Strategy}

\subsection{Design}

On every trading day, a \$1-wide bull put spread was opened on a SPY-scaled S\&P 500 (index $\div$ 10):
short a put at $-0.30$ delta, long a put \$1 lower, 21 trading days to expiry, held to expiry. Both legs
were priced by Black--Scholes using VIX as implied volatility and \texttt{\^{}IRX} as the risk-free rate.
Profit and loss was computed from the actual subsequent index level, so outcomes reflect real market
paths even though pricing is modelled.

The short strike was located by inverting the delta definition:
\begin{equation}
K_1 = S \exp\!\left( \left(r + \tfrac{\sigma^2}{2}\right)T - d_1 \sigma \sqrt{T} \right),
\qquad d_1 = -\Phi^{-1}(0.30) = 0.5244,
\end{equation}
then rounded to the nearest whole dollar to reflect actual SPY strike spacing.

Validation: realised short delta $-0.300$ against a $-0.30$ target; strikes 2.50\% below spot on
average; mean credit \$28.70; and zero violations of the theoretical P\&L bounds
$[-(100-\text{credit}), +\text{credit}]$ across all 9,201 trades.

\subsection{Performance}

\begin{table}[h]
\centering
\caption{Simulated per-trade profit and loss ($n=9{,}201$)}
\begin{tabular}{@{}lrrrrr@{}}
\toprule
Cost assumption & Mean & HAC $t$ & Subsample $t$ & Win rate & Worst \\
\midrule
Gross & \$11.85 & 9.74 & 5.98 -- 8.11 & 82.2\% & $-$\$85.41 \\
Net of \$2.00 & \$9.85 & 8.09 & 4.85 -- 6.91 & 82.0\% & $-$\$87.41 \\
Net of \$3.50 (base) & \$8.35 & 6.86 & 4.00 -- 6.02 & 82.0\% & $-$\$88.91 \\
Net of \$6.00 & \$5.85 & 4.80 & 2.59 -- 4.52 & 81.9\% & $-$\$91.41 \\
\bottomrule
\end{tabular}
\end{table}

Results are positive under every cost assumption tested. Evaluated on the second half of the sample
only (2008 onwards, $n=4{,}600$, net of \$3.50), mean profit was \textbf{\$9.72} --- \emph{stronger}
than the full sample, indicating no decay.

\subsection{Where the edge originates}

Decomposing the base case: the mean winning trade returns the full credit of \$28.70, while the mean
loss is approximately \$66. That payoff ratio implies a breakeven win rate of
\begin{equation}
\frac{66}{66 + 28.70} = 69.7\%.
\end{equation}
A fairly priced spread should therefore win roughly seven times in ten. The realised rate was
\textbf{82.2\%}. This 12.5 percentage-point gap \emph{is} the edge --- the same variance risk premium
measured in Section 3, expressed as probability rather than volatility points. It follows that a high
win rate is not itself evidence of an edge; only its excess over the breakeven rate implied by the
payoff ratio is informative.

\subsection{Volatility timing: nine filters, nine failures}

\begin{table}[h]
\centering
\caption{Filter comparison, full sample, net of \$3.50. \textit{Annual est.} $= 12 \times$ mean
$\times$ share, i.e.\ expected annual profit running one position at a time}
\begin{tabular}{@{}lrrrrrr@{}}
\toprule
Filter & $n$ & Share & Mean & Subsample $t$ & Win \% & Annual est. \\
\midrule
\textbf{No filter} & 9,201 & 1.00 & \$8.35 & 4.00 & 82.0 & \textbf{100.14} \\
VIX $\leq$ 35 & 8,848 & 0.96 & \$8.17 & 3.43 & 81.9 & 94.28 \\
VIX $\leq$ 30 & 8,467 & 0.92 & \$7.63 & 2.60 & 81.5 & 84.22 \\
VIX $\leq$ 25 & 7,603 & 0.83 & \$7.30 & 2.48 & 81.4 & 72.36 \\
VIX $\leq$ 20 & 5,778 & 0.63 & \$8.39 & 2.54 & 82.9 & 63.22 \\
VIX 15--30 & 5,520 & 0.60 & \$6.86 & 1.30 & 80.1 & 49.41 \\
VIX $\geq$ 20 & 3,427 & 0.37 & \$8.29 & 0.54 & 80.4 & 37.06 \\
VIX $\geq$ 25 & 1,601 & 0.17 & \$13.36 & 2.04 & 84.8 & 27.90 \\
VIX $\geq$ 30 & 736 & 0.08 & \textbf{\$16.53} & 1.28 & \textbf{87.8} & 15.87 \\
\bottomrule
\end{tabular}
\end{table}

No filter improves on unconditional entry. High-VIX subsets produce the highest per-trade profit
(\$16.53) and highest win rate (87.8\%) of any subset --- the opposite of the common intuition --- but
occur too rarely to build a year around. The mechanism is that delta-targeting automatically widens
the strike distance as volatility rises: a $-0.30$ delta 21-day put sits roughly 1.5\% below spot at
VIX 12 but approximately 5\% below at VIX 40, so richer premium and greater cushion arrive together.

Two features indicate the remaining differences are noise. The $\leq$ filters are non-monotonic
($\leq$20: 8.39, $\leq$25: 7.30, $\leq$30: 7.63, $\leq$35: 8.17) where a genuine effect would improve
steadily as the cap tightens. And the conservative subsample statistic falls below conventional
significance for most filters, since restricting the sample destroys the power needed to detect
anything. The out-of-sample half reproduces the same ordering.

\subsection{Annual pattern and equity beta}

Seven of thirty-seven years were unprofitable: 2000, 2001, 2002, 2007, 2008, 2018 and 2022. Every one
is an equity bear market. Since a bull put spread carries positive delta, a material share of the
simulated return is ordinary equity exposure rather than volatility premium --- a return obtainable
more cheaply from an index fund.

One contrast is instructive: 2020 was strongly profitable ($+$\$13.43 per trade) despite containing the
worst single VRP observation in the dataset, while 2022 lost \$11.31 per trade in a far less dramatic
year. A sharp crash damages a few trades and then offers rich premiums on re-entry; a prolonged decline
bleeds the strategy continuously. \emph{Duration of drawdown matters more than depth.}

\section{Discussion}

\subsection{Interpretation}

The variance risk premium exists, is large, is robust across four inference methods, and has not decayed
over 36 years. It is not a market inefficiency. The skewness of $-3.65$, the tail episodes of Table 3,
and the alignment of losing years with bear markets all indicate compensation for bearing risk that
other participants wish to shed. Sellers of volatility are paid a steady fee in exchange for absorbing
rare, severe losses --- structurally the business of an insurer.

The practical implication for strategy design is unusually simple. Every attempt to improve on mechanical,
unconditional entry made results worse. Mechanical execution is not a simplification adopted for
convenience; it is the finding.

\subsection{Implications for larger accounts}

This study measures per-trade expectancy, which is a property of the instrument rather than of the
account trading it. Capital does not alter the premium; it alters how much of the premium survives to
the account, through four channels.

\textbf{Cost drag falls with position size.} Commissions and bid--ask are approximately fixed in dollar
terms per trade, so they shrink as a proportion of a larger credit. Section 4 shows the effect
directly: costs of \$3.50 reduce mean profit from \$11.85 to \$8.35, a drag of 30\%. A wider spread
collecting proportionally more premium for the same fixed cost would retain a result closer to the
gross figure. The gap between the gross and net rows of Table 6 is therefore an upper bound on what
scale alone could recover.

\textbf{Concurrency reduces variance.} A larger account can hold several positions opened on different
dates, so performance no longer depends on which 21-day windows happened to be sampled. This does not
change expected return but should reduce the dispersion around it.

\textbf{More structures become available.} Section 4.3 shows that the long leg truncates the premium
captured: it bounds the loss but also gives back part of the credit. With sufficient collateral,
cash-secured puts would capture the untruncated premium --- at the cost of removing the floor that
bounds the outcomes in Table 3 at \$71. Given the February 2020 episode, that trade-off is not
obviously favourable at any account size.

\textbf{Sizing determines whether expectancy is realised.} A positive mean is only collected by a
participant who survives the loss distribution to collect it. Sufficient capital to size positions at a
small fraction of the account is a precondition for the results in Section 4 to translate into account
performance.

None of these channels is tested here. Wider spreads are not scaled copies of narrower ones: they carry
different net delta, different liquidity, and different assignment characteristics. The observations
above are implications of the reported results, not additional findings.

\subsection{Limitations}

\textbf{1. Volatility skew is omitted.} Both legs were priced at VIX rather than at each strike's own
implied volatility. Real out-of-the-money puts trade on a skewed surface, and VIX --- a variance-swap
construction weighting across many strikes --- is not the implied volatility of any particular option.
The sign of the resulting bias is not determinable without the surface. This is the largest unquantified
error in the study.

\textbf{2. Fills are assumed at theoretical value.} Real execution occurs at bid or ask. The cost cases
approximate this but do not model it; the difference between \$3.50 and \$6.00 of round-trip cost changes
the mean result by 30\%.

\textbf{3. Positions overlap.} Daily entry with 21-day holds implies roughly 21 concurrent positions.
Cumulative totals are therefore meaningless and only per-trade figures are reported.

\textbf{4. HAC and subsample statistics diverge} throughout Section 4, unlike Section 3 where they agreed.
The assumption-free subsample range is the more credible figure.

\textbf{5. Return on risk capital remains high} relative to CBOE's published put-writing benchmarks. Part
is explained by the structural leverage of defined-risk spreads, but the gap is not fully reconciled.

\textbf{6. Equity beta is not separated} from volatility premium.

\textbf{7. No trade management.} Positions were held to expiry, unlike practical implementations that exit
at partial profit.

\textbf{8. Position sizing is not addressed.} The worst simulated trade lost \$88.91. Whether a
sequence of such losses is survivable depends on account size and sizing policy, neither of which is
examined here. Expectancy and survival are separate questions; this study addresses only the first.

\section{Conclusion}

Implied volatility on S\&P 500 index options has exceeded subsequent realised volatility by an average of
4.10 volatility points since 1990. The effect is statistically robust under correct treatment of
overlapping windows, present in every decade examined, and captured in simulation by a defined-risk
option spread returning \$8.35 per trade net of costs.

The premium is compensation for crash risk rather than a market inefficiency, and its risk-adjusted
attractiveness has deteriorated markedly since the 1990s even as its magnitude has held. Attempts to
time entry by volatility regime uniformly failed.

Whether the effect is exploitable in practice remains open, and turns on two assumptions this study
could not test: the shape of the volatility surface at the strikes actually traded, and the cost of
real execution. Both require option chain data rather than index data. That extension --- repricing the
same trades against historical quotes from OptionMetrics --- is the natural continuation of this work.

\section*{Reproducibility}

All analysis was conducted in Python (pandas, NumPy, SciPy, statsmodels) using publicly available data.
Code, methodology specifications with pre-registered falsification criteria, and full result outputs are
available on request. Falsification criteria were committed to in writing before data examination in both
studies.

\section*{Note on tools}

Analysis code was written by the author. Methodology design, statistical review and manuscript
preparation were carried out in collaboration with an AI assistant (Claude, Anthropic). All results were
computed and verified by the author.

\begin{thebibliography}{9}
\small
\bibitem{bs} Black, F. and Scholes, M. (1973). The Pricing of Options and Corporate Liabilities.
\textit{Journal of Political Economy}, 81(3), 637--654.
\bibitem{bakshi} Bakshi, G. and Kapadia, N. (2003). Delta-Hedged Gains and the Negative Market Volatility
Risk Premium. \textit{Review of Financial Studies}, 16(2), 527--566.
\bibitem{carrwu} Carr, P. and Wu, L. (2009). Variance Risk Premiums. \textit{Review of Financial Studies},
22(3), 1311--1341.
\bibitem{btz} Bollerslev, T., Tauchen, G. and Zhou, H. (2009). Expected Stock Returns and Variance Risk
Premia. \textit{Review of Financial Studies}, 22(11), 4463--4492.
\bibitem{nw} Newey, W. K. and West, K. D. (1987). A Simple, Positive Semi-Definite, Heteroskedasticity and
Autocorrelation Consistent Covariance Matrix. \textit{Econometrica}, 55(3), 703--708.
\bibitem{cboe} Chicago Board Options Exchange (2019). \textit{Cboe Volatility Index (VIX) White Paper}.
\end{thebibliography}

\end{document}
